\runninghead{The Appointed Texts}
\section{The Appointed Texts}
\label{sec:appointed-texts}

The Latin below is that of the 1962 typical edition, \work{Missale Romanum}
(Typis Polyglottis Vaticanis), printed pp.~422--424, nos.~1769--1778, read on
the page images of the CMAA facsimile. The Mass is the last of the Proper of
Time, printed whole under its own heading, \latin{Dominica vigesima quarta et
ultima post Pentecosten}, with its five chants written out. Every word of the
ten texts stands in the Ratisbon printing of the \work{Missale Romanum}
(Pustet, 1862), pp.~359--360, under the heading \latin{Dominica XXIV.\ et
ultima post Pentecosten}. The two printings differ in spelling
(\latin{exequéntes}, \latin{intélligat}, \latin{cœlo}), in pointing,
capitals and the form of the citations, and in the Secret's conclusion cue,
and in no word that is said.

The English of the three orations and of the five chants is that of the
Lasance \work{New Roman Missal} in the Benziger revision of 1945, pp.~736--740,
which prints this Mass in Latin and English under the heading ``Twenty-fourth
and Last Sunday after Pentecost.'' Its English of the chants renders the
Missal's own wording where that departs from the Clementine Vulgate. The
Douay--Rheims Bible in Bishop Challoner's revision, which translates the
Clementine Vulgate, is given for the two lessons, for the Introit's psalm
verse, and beside each other chant at the Bible verses it comes from; each
place where the Missal's wording departs from the Vulgate is named beneath
the text. None of this English is an approved liturgical translation of the
1962 Missal. Psalms are cited in the Missal's numbering, with the Hebrew
number in parentheses at each psalm text below.

\subsection{Introit}
\label{proper-introit}

\properlocus{\latin{Antiphona ad Introitum}, \latin{Ierem.\ 29, 11, 12 et 14};
verse \latin{Ps.\ 84, 2} (Heb.\ 85). Missal p.~422, no.~1769.}

\begin{latinproper}
Dicit Dóminus: Ego cógito cogitatiónes pacis, et non afflictiónis:
invocábitis me, et ego exáudiam vos: et redúcam captivitátem vestram de
cunctis locis. \textnormal{Ps.\ 84, 2} Benedixísti, Dómine, terram tuam:
avertísti captivitátem Iacob. ℣.~Glória Patri.
\end{latinproper}

\begin{englishwitness}{Lasance, The New Roman Missal (1945), p.~736, Introit}
The Lord saith, I think thoughts of peace, and not of affliction; you shall
call upon Me and I will hear you; and I will bring back your captivity from
all places. Ps.~84, 2. Lord, Thou hast blessed Thy land: Thou hast turned away
the captivity of Jacob. ℣.~Glory.
\end{englishwitness}

\begin{englishwitness}{Douay--Rheims (Challoner), Jer 29:11--12, 14, the antiphon's verses}
For I know the thoughts that I think towards you, saith the Lord, thoughts of
peace, and not of affliction, to give you an end and patience. And you shall
call upon me, and you shall go: and you shall pray to me, and I will hear you.
\dots\ And I will be found by you, saith the Lord: and I will bring back your
captivity, and I will gather you out of all nations, and from all the places
to which I have driven you out, saith the Lord: and I will bring you back from
the place to which I caused you to be carried away captive.
\end{englishwitness}

\begin{englishwitness}{Douay--Rheims (Challoner), Ps 84:2, the verse}
Lord, thou hast blessed thy land: thou hast turned away the captivity of
Jacob.
\end{englishwitness}

\noindent The antiphon is not a continuous text of the Vulgate. It opens
\latin{Dicit Dóminus}, ``The Lord saith,'' where the Vulgate's verse~11 reads
\latin{Ego enim scio cogitationes quas ego cogito super vos, ait Dominus}. It
has nothing of \latin{ut dem vobis finem et patientiam} (``to give you an end
and patience''), of \latin{et ibitis: et orabitis me} (``and you shall go: and
you shall pray to me''), of verse~13, of \latin{Et inveniar a vobis, ait
Dominus} (``And I will be found by you, saith the Lord''), or of \latin{et
congregabo vos de universis gentibus} (``and I will gather you out of all
nations''). It joins \latin{redúcam captivitátem vestram} directly to
\latin{de cunctis locis}, so that the Lord promises to bring the captivity
back ``from all places.'' The psalm verse agrees with the Vulgate. After the
\latin{Gloria Patri} the antiphon is repeated.

\subsection{Collect}
\label{proper-collect}

\properlocus{\latin{Oratio}. Missal p.~423, no.~1770.}

\begin{latinproper}
Excita, quǽsumus, Dómine, tuórum fidélium voluntátes: ut, divíni óperis
fructum propénsius exsequéntes; pietátis tuæ remédia maióra percípiant. Per
Dóminum.
\end{latinproper}

\begin{englishwitness}{Lasance, The New Roman Missal (1945), p.~737, Prayer}
Arouse, we beseech Thee, O Lord, the wills of Thy faithful, that, by more
earnestly following after the fruit of the divine work, they may the more
abundantly partake of Thy mercies. Through our Lord.
\end{englishwitness}

\noindent The Lasance English is free at the end. Where it has ``the more
abundantly partake of Thy mercies,'' the Latin's object is \latin{remédia},
``remedies,'' and the comparative is \latin{maióra}, ``greater.''

\subsection{Epistle}
\label{proper-epistle}

\properlocus{\latin{Léctio Epístolæ beáti Pauli Apóstoli ad Colossénses},
Col 1:9--14. Missal p.~423, no.~1771.}

\begin{latinproper}
Fratres: Non cessámus pro vobis orántes, et postulántes ut impleámini
agnitióne voluntátis Dei, in omni sapiéntia et intelléctu spiritáli: ut
ambulétis digne Deo per ómnia placéntes: in omni ópere bono fructificántes,
et crescéntes in sciéntia Dei: in omni virtúte confortáti secúndum poténtiam
claritátis eius in omni patiéntia, et longanimitáte cum gáudio, grátias
agéntes Deo Patri, qui dignos nos fecit in partem sortis sanctórum in lúmine:
qui erípuit nos de potestáte tenebrárum, et tránstulit in regnum Fílii
dilectiónis suæ, in quo habémus redemptiónem per sánguinem eius, remissiónem
peccatórum.
\end{latinproper}

\begin{englishwitness}{Douay--Rheims (Challoner), Col 1:9--14}
Therefore we also, from the day that we heard it, cease not to pray for you
and to beg that you may be filled with the knowledge of his will, in all
wisdom and spiritual understanding: That you may walk worthy of God, in all
things pleasing; being fruitful in every good work and increasing in the
knowledge of God: Strengthened with all might according to the power of his
glory, in all patience and longsuffering with joy, Giving thanks to God the
Father, who hath made us worthy to be partakers of the lot of the saints in
light: Who hath delivered us from the power of darkness and hath translated us
into the kingdom of the Son of his love, In whom we have redemption through
his blood, the remission of sins.
\end{englishwitness}

\noindent The lesson begins within verse~9. The Missal sets \latin{Fratres:
Non cessámus}, which the Lasance renders ``Brethren, We cease not,'' where the
Vulgate reads \latin{Ideo et nos ex qua die audivimus, non cessamus}
(``Therefore we also, from the day that we heard it, cease not''). In the same
verse the Missal reads \latin{agnitióne voluntátis Dei}, ``the knowledge of
the will of God,'' as the Lasance has it, where the Vulgate reads
\latin{voluntatis ejus}, ``his will.'' Every other word agrees with the
Vulgate. The Missal ends the lesson at verse~14 with a full stop; in the
Vulgate the sentence runs on into the praise of the Son, ``Who is the image of
the invisible God'' (v.~15).

\subsection{Gradual}
\label{proper-gradual}

\properlocus{\latin{Graduale}, \latin{Ps.\ 43, 8-9} (Heb.\ 44). Missal p.~423,
no.~1772.}

\begin{latinproper}
Liberásti nos, Dómine, ex affligéntibus nos: et eos, qui nos odérunt,
confudísti. ℣.~In Deo laudábimur tota die, et in nómine tuo confitébimur in
sǽcula.
\end{latinproper}

\begin{englishwitness}{Lasance, The New Roman Missal (1945), pp.~737--738, Gradual}
Thou hast delivered us, O Lord, from them that afflict us: and hast put them
to shame that hate us. ℣.~In God we will glory all the day: and in Thy name we
will give praise for ever.
\end{englishwitness}

\begin{englishwitness}{Douay--Rheims (Challoner), Ps 43:8--9}
But thou hast saved us from them that afflict us: and hast put them to shame
that hate us. In God shall we glory all the day long: and in thy name we will
give praise for ever.
\end{englishwitness}

\noindent The Vulgate's verse~8 reads \latin{salvasti enim nos de affligentibus
nos, et odientes nos confudisti}. The Gradual has \latin{Liberásti},
``thou hast delivered,'' for \latin{salvasti}, ``thou hast saved''; adds
\latin{Dómine}; has \latin{ex} for \latin{de}; and has \latin{eos, qui nos
odérunt}, ``those who hated us,'' for \latin{odientes nos}, ``those hating
us.'' Verse~9 ends \latin{in saeculum} in the Vulgate and \latin{in sǽcula}
in the Gradual.

\subsection{Alleluia}
\label{proper-alleluia}

\properlocus{\latin{Ps.\ 129, 1-2} (Heb.\ 130). Missal p.~423, no.~1773.}

\begin{latinproper}
Allelúia, allelúia. ℣.~De profúndis clamávi ad te, Dómine: Dómine, exáudi
oratiónem meam. Allelúia.
\end{latinproper}

\begin{englishwitness}{Lasance, The New Roman Missal (1945), p.~738, Alleluia}
Alleluia, alleluia. ℣.~From the depths I have cried to Thee, O Lord: Lord,
hear my prayer. Alleluia.
\end{englishwitness}

\begin{englishwitness}{Douay--Rheims (Challoner), Ps 129:1--2}
Out of the depths I have cried to thee, O Lord: Lord, hear my voice. Let thy
ears be attentive to the voice of my supplication.
\end{englishwitness}

\noindent The Vulgate reads \latin{Domine, exaudi vocem meam}, ``Lord, hear my
voice''; the chant reads \latin{oratiónem meam}, ``my prayer,'' and stops
before the second half of verse~2. The Mass has no Tract and no Sequence.

\subsection{Gospel}
\label{proper-gospel}

\properlocus{\latin{Sequéntia sancti Evangélii secúndum Matthǽum}, Mt
24:15--35. Missal pp.~423--424, no.~1774; the printed cue \latin{Credo}
follows.}

\begin{latinproper}
In illo témpore: Dixit Iesus discípulis suis: Cum vidéritis abominatiónem
desolatiónis, quæ dicta est a Daniéle prophéta, stantem in loco sancto: qui
legit, intéllegat: tunc qui in Iudǽa sunt, fúgiant ad montes: et qui in
tecto, non descéndat tóllere áliquid de domo sua: et qui in agro, non
revertátur tóllere túnicam suam. Væ autem prægnántibus et nutriéntibus in
illis diébus. Oráte autem, ut non fiat fuga vestra in híeme vel sábbato. Erit
enim tunc tribulátio magna, qualis non fuit ab inítio mundi usque modo,
neque fiet. Et nisi breviáti fuíssent dies illi, non fíeret salva omnis
caro: sed propter eléctos breviabúntur dies illi. Tunc si quis vobis
díxerit: Ecce hic est Christus, aut illic: nolíte crédere. Surgent enim
pseudochrísti, et pseudoprophétæ: et dabunt signa magna et prodígia, ita ut
in errórem inducántur (si fíeri potest) étiam elécti. Ecce prædíxi vobis. Si
ergo díxerint vobis: Ecce in desérto est, nolíte exíre: ecce in
penetrálibus, nolíte crédere. Sicut enim fulgur exit ab Oriénte, et paret
usque in Occidéntem: ita erit et advéntus Fílii hóminis. Ubicúmque fúerit
corpus, illic congregabúntur et áquilæ. Statim autem post tribulatiónem
diérum illórum sol obscurábitur, et luna non dabit lumen suum, et stellæ
cadent de cælo, et virtútes cælórum commovebúntur: et tunc parébit signum
Fílii hóminis in cælo: et tunc plangent omnes tribus terræ: et vidébunt
Fílium hóminis veniéntem in núbibus cæli cum virtúte multa, et maiestáte. Et
mittet Angelos suos cum tuba et voce magna: et congregábunt eléctos eius a
quátuor ventis, a summis cælórum usque ad términos eórum. Ab árbore autem
fici díscite parábolam: cum iam ramus eius tener fúerit, et fólia nata,
scitis quia prope est æstas: ita et vos cum vidéritis hæc ómnia, scitóte
quia prope est in iánuis. Amen dico vobis, quia non præteríbit generátio
hæc, donec ómnia hæc fiant. Cælum et terra transíbunt, verba autem mea non
præteríbunt.
\end{latinproper}

\begin{englishwitness}{Douay--Rheims (Challoner), Mt 24:15--35}
When therefore you shall see the abomination of desolation, which was spoken
of by Daniel the prophet, standing in the holy place: he that readeth let him
understand. Then they that are in Judea, let them flee to the mountains: And
he that is on the housetop, let him not come down to take any thing out of his
house: And he that is in the field, let him not go back to take his coat. And
woe to them that are with child and that give suck in those days. But pray
that your flight be not in the winter or on the sabbath. For there shall be
then great tribulation, such as hath not been from the beginning of the world
until now, neither shall be. And unless those days had been shortened, no
flesh should be saved: but for the sake of the elect those days shall be
shortened. Then if any man shall say to you, Lo here is Christ, or there: do
not believe him. For there shall arise false Christs and false prophets and
shall shew great signs and wonders, insomuch as to deceive (if possible) even
the elect. Behold I have told it to you, beforehand. If therefore they shall
say to you, Behold he is in the desert: go ye not out. Behold he is in the
closets: believe it not. For as lightning cometh out of the east and appeareth
even into the west: so shall also the coming of the Son of man be. Wheresoever
the body shall be, there shall the eagles also be gathered together. And
immediately after the tribulation of those days, the sun shall be darkened
and the moon shall not give her light and the stars shall fall from heaven and
the powers of heaven shall be moved. And then shall appear the sign of the Son
of man in heaven. And then shall all tribes of the earth mourn: and they shall
see the Son of man coming in the clouds of heaven with much power and majesty.
And he shall send his angels with a trumpet and a great voice: and they shall
gather together his elect from the four winds, from the farthest parts of the
heavens to the utmost bounds of them. And from the fig tree learn a parable:
When the branch thereof is now tender and the leaves come forth, you know that
summer is nigh. So you also, when you shall see all these things, know ye that
it is nigh, even at the doors. Amen I say to you that this generation shall
not pass till all these things be done. Heaven and earth shall pass: but my
words shall not pass.
\end{englishwitness}

\noindent The Missal's opening, \latin{In illo témpore: Dixit Iesus discípulis
suis: Cum vidéritis}, which the Lasance renders ``At that time, Jesus said to
His disciples, When you shall see,'' stands for the Vulgate's \latin{Cum ergo
videritis}, ``When therefore you shall see.'' The Missal drops the
\latin{ergo}, ``therefore,'' which in Matthew ties the saying to the
disciples' question on the Mount of Olives (Mt 24:3). It spells
\latin{intéllegat} where the Vulgate has \latin{intelligat}. Every other word
of verses 15--35 agrees with the Vulgate.

\subsection{Offertory}
\label{proper-offertory}

\properlocus{\latin{Antiphona ad Offertorium}, \latin{Ps.\ 129, 1-2} (Heb.\
130). Missal p.~424, no.~1775.}

\begin{latinproper}
De profúndis clamávi ad te, Dómine: Dómine, exáudi oratiónem meam: de
profúndis clamávi ad te, Dómine.
\end{latinproper}

\begin{englishwitness}{Lasance, The New Roman Missal (1945), p.~739, Offertory}
From the depths I have cried out to Thee, O Lord: Lord, hear my prayer: from
the depths I have cried out to Thee, O Lord.
\end{englishwitness}

\noindent The Offertory sings the Alleluia's verses again, with the same
\latin{oratiónem meam}, ``my prayer,'' where the Vulgate has ``my voice''
(the Douay text stands under the Alleluia above), and repeats its first line
at the end.

\subsection{Secret}
\label{proper-secret}

\properlocus{\latin{Secreta}. Missal p.~424, no.~1776; the printed direction
\latin{Præfatio de Ssma Trinitate} follows.}

\begin{latinproper}
Propítius esto, Dómine, supplicatiónibus nostris: et, pópuli tui oblatiónibus
precibúsque suscéptis, ómnium nostrum ad te corda convérte; ut, a terrénis
cupiditátibus liberáti, ad cæléstia desidéria transeámus. Per Dóminum
nostrum.
\end{latinproper}

\begin{englishwitness}{Lasance, The New Roman Missal (1945), p.~740, Secret}
Be propitious, O Lord, to our supplications, and accept the offerings and
prayers of Thy people: turn all our hearts unto Thee, that, being delivered
from earthly desires, we may pass on to the enjoyments of heaven. Through our
Lord.
\end{englishwitness}

\noindent Lasance's ``earthly desires'' renders \latin{terrénis
cupiditátibus}, and his ``the enjoyments of heaven'' stands for \latin{cæléstia
desidéria}, in which \latin{desidéria} is ``desires'': the Latin names desire
on both sides of the passage.

\subsection{Communion}
\label{proper-communion}

\properlocus{\latin{Antiphona ad Communionem}, \latin{Marc.\ 11, 24}. Missal
p.~424, no.~1777.}

\begin{latinproper}
Amen dico vobis, quidquid orántes pétitis, crédite quia accipiétis, et fiet
vobis.
\end{latinproper}

\begin{englishwitness}{Lasance, The New Roman Missal (1945), p.~740, Communion}
Amen, I say to you, whatsoever you ask when you pray, believe that you shall
receive, and it shall be done unto you.
\end{englishwitness}

\begin{englishwitness}{Douay--Rheims (Challoner), Mk 11:24}
Therefore I say unto you, all things, whatsoever you ask when ye pray, believe
that you shall receive: and they shall come unto you.
\end{englishwitness}

\noindent The Vulgate's verse reads \latin{Propterea dico vobis, omnia
quaecumque orantes petitis, credite quia accipietis, et evenient vobis}. The
antiphon differs in three places: it opens \latin{Amen dico vobis}, which
Lasance renders ``Amen, I say to you''; it has \latin{quidquid},
``whatsoever,'' for \latin{omnia quaecumque}, ``all things, whatsoever''; and
it closes \latin{et fiet vobis}, ``and it shall be done unto you.''

\subsection{Postcommunion}
\label{proper-postcommunion}

\properlocus{\latin{Postcommunio}. Missal p.~424, no.~1778.}

\begin{latinproper}
Concéde nobis, quǽsumus, Dómine: ut per hæc sacraménta quæ súmpsimus,
quidquid in nostra mente vitiósum est, ipsórum medicatiónis dono curétur. Per
Dóminum.
\end{latinproper}

\begin{englishwitness}{Lasance, The New Roman Missal (1945), p.~740, Postcommunion}
Grant us, we beseech Thee, O Lord, that, through this sacrament which we have
received, whatever is evil in our hearts may be restored by its gift of
healing. Through our Lord.
\end{englishwitness}

\noindent The Latin differs from the Lasance English in three words.
\latin{Sacraménta} is plural, ``sacraments''; \latin{mente} is ``mind,'' where
Lasance has ``hearts''; and \latin{vitiósum} is ``faulty,'' where Lasance has
``evil.''
